% ============================================
% Section: Discussion
% ============================================
\label{sec:discussion}

\subsection{Design Decisions}

\paragraph{Why an application-layer protocol?} Off-the-shelf transports
(ROS~2/DDS with default QoS, gRPC, raw TCP) provide no differentiation
between a 20~Hz pose packet and a 3~MB submap upload, resulting in
head-of-line blocking, no deadline awareness, and no compression integration.
Co-LIC2's protocol is a thin layer ($\sim$2000~LoC of C++) that adds
exactly the QoS differentiation needed for cooperative SLAM and deploys
alongside existing middleware rather than replacing it.

\paragraph{Why cluster-head election versus static topology?} Static
topologies fail in mobile deployments: a pre-designated head may move
behind a building, losing 50\% of its members. Our election protocol adapts
on a $\sim$3~s timescale, comparable to submap publication rate ($\sim$3~s)
and merge latency ($<$2~s).

\paragraph{Why NPV as the compression objective?} Perceptual metrics (PSNR,
SSIM, LPIPS) require rendering---too expensive ($\sim$15~ms at 720p) for
per-submap optimization on an edge GPU already occupied with real-time
SLAM. NPV computes in 1~ms from Gaussian parameters alone and empirically
correlates with rendering quality ($\rho\approx 0.89$ in outdoor scenes).

\subsection{Limitations}

\paragraph{Static-scene assumption.} The per-node front-end inherits
Gaussian-LIC2's static-scene assumption. Confidence-weighted conflict
resolution handles slowly changing scenes (parked cars, construction
progress) but not real-time dynamic-object tracking. Integration with
RoSe-SLAM~\cite{wang2026roseslam}-style motion segmentation is the most
important next step.

\paragraph{EDCA dependence.} Our QoS mapping assumes 802.11e EDCA, which
requires WMM support on the AP. In IBSS mode EDCA is optional; in 5G~NR
deployments per-flow QoS requires network-slice configuration. Without
EDCA, all streams fall to AC\_BE and latency bounds degrade
proportionally.

\paragraph{No encryption.} The protocol authenticates via signed
descriptors and submaps (Ed25519) but does not encrypt traffic. Trusted
environments need TLS or application-layer encryption over submap and
keyframe streams.

\paragraph{Head scalability ceiling.} A single head handles $\sim$8 members
at $\sim$17~Mbps ingress. Cross-cluster consistency depends on the Tier-0
server or opportunistic head-to-head exchange. Practical limit: $\sim$32
nodes with 4 heads + 1 server.

\subsection{Open Questions}

\begin{itemize}
  \item \textbf{Cluster splitting/merging.} Should an over-subscribed head
    split the cluster, or degrade QoS? A split protocol is complex but
    preserves latency bounds.
  \item \textbf{Energy-optimal UAV compression.} For battery-constrained
    UAVs, the trade-off between local compression (GPU energy) and
    transmission (radio energy) is deployment-specific.
  \item \textbf{5G~NR sliced-network campaign.} Our 5QI mapping is based on
    3GPP specifications; a measurement campaign on a real 5G~SA network
    would validate the assumptions.
  \item \textbf{Multi-node dynamic consensus.} Per-node dynamic filtering
    plus cross-node voting to agree on the static subset of the map for
    long-term urban deployment.
\end{itemize}
